%!TeX root=../thesis.tex

\chapter{Conclusion}
\label{ch:conclusion}
\label{chap:conclusion}

This final chapter synthesizes what was learned and what it implies for practice and for further research.
The empirical results and the refinements to Hypo-Stage and its usage guidelines are reported together in Chapter~\ref{chap:results}: Sections~\ref{sec:results-rq1}--\ref{sec:results-rq4} present the findings tied to RQ1--RQ4, and Section~\ref{sec:refinements} documents the tool and guideline changes driven by the field study (especially RQ2).
That integrated layout matches the formative, iterative research design described in Chapter~\ref{chap:research-design}: observation and refinement form one action-research loop rather than a detached annex.

\section{Summary in Relation to the Research Questions}
\label{sec:conclusion-rq-summary}

% TODO: Answer RQ1--RQ4 explicitly, referencing evidence from Chapter~\ref{chap:results}.

\section{Implications for Practice and Research}
\label{sec:conclusion-implications}

% TODO: Practitioners (Hypo-Stage, ArchHypo, IDPs); researchers (tool-mediated architectural hypotheses, case-study generalization).

\section{Limitations and Future Work}
\label{sec:conclusion-limitations}

% TODO: Organizational scope; Backstage coupling; length of observation; validation of refinements in further cycles.
